\section{みなし相続財産}

\subsection{相続税の計算}

\begin{shortanswer}[年金1 2022 問題2(5)ア]
  相続人が取得した退職手当金等については、課税価格の計算において一定の非課税限度額が
あるが、その内容について簡潔に入力しなさい。
\end{shortanswer}

\begin{answerbox}
  $法定相続人の数 \times 500万円$が非課税限度額となる。
\end{answerbox}


\begin{shortanswer}[年金1 2022 問題2(5)イ]
  相続税の計算では、みなし相続財産等を含む課税対象財産から計算した課税価格に対し、基
礎控除額があるが、その内容について簡潔に入力しなさい。
\end{shortanswer}

\begin{answerbox}
  $3,000万円 + 法定相続人の数 \times 600万円$が基礎控除額となる。
\end{answerbox}

\newpage 

\subsection{その他}

\begin{truefalse}[年金1 H29 問題1(6)\textcircled{5}]
           \textbf{下線部の正誤を判定し、誤っている場合は正しい内容に改めなさい。}
\\  

  確定給付企業年金の遺族給付金は、被相続人の死亡後\underline{ 6 か月以内}に支給額が確定したものに
限り、相続税法上みなし相続財産（退職手当金等に含まれる給付）として相続税の課税対象と
なる。
\end{truefalse}

\begin{answerbox}
  \textbf{誤}

  3年以内
\end{answerbox}

\newpage